\documentclass[11pt,a4paper]{article}
\usepackage{notes}
\usepackage{fancyhdr}

\pagestyle{fancy}
\fancyhf{}
\lhead{Geophysics IIIA: Potential Fields \& Geothermics}
\rhead{Week 8 Thermal Refraction -- 2}
\cfoot{\thepage}

\title{Week 8 --- Activity: Thermal Refraction in Layered Crustal Rocks}
\author{Geophysics 3A: Potential Fields \& Geothermics}
\date{\today}

\begin{document}
\maketitle

\section*{Purpose}
To explore thermal refraction caused by conductivity contrasts within the crust, particularly in sedimentary and metamorphic layering.

\section*{Geological Context}
Many crustal sections contain cyclic layering inherited from sedimentary
environments, such as alternating shales and quartz-rich sandstones that later
become quartzites.

\begin{itemize}
  \item Shales: low thermal conductivity.
  \item Quartzites: high thermal conductivity.
\end{itemize}

This layering can persist through burial, deformation, and metamorphism.

\section*{Tasks}
\begin{enumerate}
  \item Consider a vertical sequence of alternating shale and quartzite layers.
  Assume constant heat flow through the sequence.

  \item Using
  \[
    q = -k \nabla T,
  \]
  explain how the temperature gradient differs in shale layers compared to quartzite layers.

  \item Sketch a qualitative temperature--depth profile showing changes in slope across layers.

  \item Now consider the same layered sequence tilted relative to vertical heat flow. Describe how thermal refraction will distort isotherms.
\end{enumerate}

\section*{Extension}
\begin{itemize}
  \item How can cyclic layering create an effectively anisotropic thermal conductivity?
  \item How might folding or tilting further modify the thermal field?
\end{itemize}

\section*{Connection}
Relate this behavior to the physical-properties mixing demonstrations discussed earlier in the course.

\end{document}